\begin{abstract}
In traditional data-driven fault diagnosis, acquiring sufficient training examples for all potential fault classes is highly challenging, particularly for rare or unseen conditions. While existing zero-shot learning methods based on Generative Adversarial Networks (GANs) have attempted to address this, they often suffer from training instability, mode collapse, and a lack of robustness in complex industrial environments. To overcome these limitations, this paper proposes a novel CLIP-Guided Dual-Path Diffusion Model (CDDM) for robust zero-shot fault diagnosis. Specifically, the CDDM leverages Large Language Models (LLMs) to generate continuous semantic embeddings, which are then aligned across modalities using a CLIP-based framework. A dual-path diffusion mechanism is introduced to capture features in both the time and frequency domains simultaneously, while a Cross-Attention mechanism effectively injects semantic information into the denoising process. Furthermore, we implement a Multi-Modal Contrastive Consistency Loss to ensure strict alignment between synthetic features and their corresponding semantic descriptions. Extensive experiments conducted on three representative datasets—TPTS, TEP, and the Hydraulic System—demonstrate that the proposed CDDM significantly outperforms state-of-the-art methods, achieving a 12.3\% improvement in diagnostic accuracy. Our results highlight the potential of diffusion-based cross-modal alignment for identifying unseen faults in intelligent manufacturing systems.
\end{abstract}

\begin{IEEEkeywords}
Zero-shot learning, fault diagnosis, dual-path diffusion model, CLIP-guided alignment, cross-modal learning.
\end{IEEEkeywords}
